\appendix

%==============================================================================
\chapter{Mathematical Derivations \& Loss Formulations}
\label{app:math}

\draftnote{Expand each of the following from the definitions already given in
  Chapter~\ref{ch:methodology}, with full derivations rather than statements.}

\begin{itemize}
  \item \textbf{Cosine reconstruction loss} for the GAAE, with both the topological and
    feature terms and the weighting between them.
  \item \textbf{VGAE evidence lower bound} --- derivation of the KL term against the
    standard normal prior, and the posterior-collapse condition that motivates the
    anti-collapse configurations.
  \item \textbf{Weighted binary cross-entropy} with $\texttt{pos\_weight} =
    n_{\mathrm{neg}} / n_{\mathrm{pos}}$, and its relationship to threshold placement.
  \item \textbf{Youden's $J$ and the $F_1$-optimal point} as threshold rules, and why both
    must be estimated out-of-fold (Section~\ref{sec:thresholds}).
  \item \textbf{Cox partial likelihood}, retained for the scoped survival extension
    (Section~\ref{sec:prognoser-scope}).
\end{itemize}

%==============================================================================
\chapter{Data Split Integrity \& Overlap Verification Logs}
\label{app:splits}

\draftnote{This appendix must be \emph{evidence}, not summary. Reproduce the actual
  \texttt{run\_full\_audit} output for the downstream splits --- the pairwise-disjoint
  assertions and per-split subject counts --- as emitted by the
  \texttt{sanity-split-hygiene} experiment, plus the cohort-parity verification output for
  the matched $2 \le T \le 3$ window showing byte-level subject-identifier and fold-membership
  agreement between the two models compared in Chapter~\ref{ch:sota-repair}. Generate,
  do not transcribe.}

%==============================================================================
\chapter{Hyperparameters \& Complete Cross-Validation Metrics}
\label{app:hyperparams}

\draftnote{Auto-generate from the results ledger (\texttt{CLASSIFIER/outputs/RESULTS.csv})
  and the per-run \texttt{resolved\_config.json} artifacts described in
  Section~\ref{sec:runner}. Two tables: (i) the resolved hyperparameters of every
  experiment referenced in Chapters~\ref{ch:sota-repair} and~\ref{ch:results}, and (ii) the
  complete per-fold cross-validation metrics behind every summary statistic quoted in those
  chapters. Hand-transcription is explicitly not acceptable here --- the ledger exists so
  that this appendix cannot drift from the runs.}

%==============================================================================
\chapter{Software Dependencies \& Hardware Environment}
\label{app:environment}

\begin{table}[htbp]
  \centering
  \small
  \begin{tabular}{@{}ll@{}}
    \toprule
    \textbf{Component} & \textbf{Version} \\
    \midrule
    Python              & 3.10.12 \\
    PyTorch             & 2.10.0+cu128 \\
    PyTorch Geometric   & 2.7.0 \\
    NumPy               & 1.26.4 \\
    pandas              & 2.2.0 \\
    scikit-learn        & 1.7.2 \\
    \midrule
    GPU                 & NVIDIA TITAN RTX, 24\,GB \\
    Driver              & 550.163.01 \\
    Hosts               & two identical workstations on a shared network export \\
    \bottomrule
  \end{tabular}
  \caption{Pinned software and hardware environment for every run reported in this thesis.}
  \label{tab:environment}
\end{table}

\draftnote{Add the fMRIPrep and denoising toolchain versions from the preprocessing cluster,
  which are not part of the project virtual environment.}

\paragraph{Determinism settings, stated explicitly.} All runs set the standard random seeds
and \texttt{cudnn.deterministic}. No run sets
\texttt{torch.use\_deterministic\_algorithms(True)} or \texttt{CUBLAS\_WORKSPACE\_CONFIG}.
This is disclosed rather than omitted because it is the exact condition under which the
non-determinism of Section~\ref{sec:btgt-reproducibility} arises, and a reader needs it to
interpret every error bar in this thesis.

%==============================================================================
\chapter{Reviewer Risk Register}
\label{app:register}

Every annotation raised in the body of this thesis is collected below, in order of
appearance, with its page. Four kinds appear:

\begin{description}[leftmargin=3.2cm,style=nextline]
  \item[RESEARCH GAP] An open problem in the field that this work engages with.
  \item[GREY AREA] A point on which the evidence is genuinely contested.
  \item[OUR EDGE] A point where this work is defensibly ahead of published practice.
  \item[EXPOSURE] A point where this work is vulnerable --- what a hostile reviewer, or an
    examiner, should be expected to attack first.
\end{description}

The register is assembled by the typesetting engine from the annotations themselves, so it
cannot fall out of sync with the text. The exposures are listed as prominently as the
strengths deliberately: a defence is easier when the weaknesses have already been named by
the candidate.

\listofannotations
